\chapter{State of the Art}
\label{chap:state-of-the-art}

\section{Social cognition, interactional legibility, and non-verbal design}
\label{sec:soa-legibility}

Occupying the same room as a person does not make a robot socially embedded. \citet{Dautenhahn2002SociallyEmbedded} define social embeddedness through structural coupling with a social environment: the agent's sensorimotor activity is shaped by an environment that contains other agents, and in turn changes it. Social signals added to a controller that otherwise ignores other agents do not meet this definition, so the design problem lies inside the controller. \citet{Henschel2020SocialCognition} frame HRI in the same spirit, as a setting in which people bring their own social-cognitive mechanisms to bear on what a robot does, and \citet{Sandini2025MutualUnderstanding} add that mutual understanding rests on shared embodied cognition and accumulated first-person experience. The engineering requirement that follows does not assume the robot has human social cognition: its behaviour must let people infer what it is attending to, what it is trying to do, and how their own behaviour changes it.

Gaze is the channel where this requirement has been tested most. \citet{Admoni2017SocialEyeGaze} separate the literature into human-focused, design-focused, and technology-focused traditions. The split matters because the same head movement is a perceptual stimulus to one tradition, a design choice to another and controller output to a third. Mutual gaze signals availability and referential gaze establishes joint attention, while gaze aversion regulates conversational timing. Whether these functions transfer to a given robot depends on its body. Robots differ in eye articulation, head--eye coordination, field of view and how visible their gaze direction is, so findings from expressive virtual agents or from one eye mechanism do not automatically hold for another.

In a controlled hallway study, \citet{Fiore2013SocialCues} found that a robot's proxemic behaviour, but not its gaze, altered perceived social presence and emotional attribution: a robot that slowed and yielded scored higher than one that sped up and cut in front, and attributions grew over repeated passes whatever the proxemic condition. \citet{He2024RobotGaze} showed that gaze during navigation can also shape perceived social presence, but the preferred gaze depends on the scenario: a glance at the person is preferred where paths cross and adds little where they do not. \citeauthor{Saunderson2019Nonverbal}'s (\citeyear{Saunderson2019Nonverbal}) survey reaches a similar conclusion across body movement (kinesics), use of space (proxemics), touch (haptics) and timing (chronemics). Non-verbal behaviour affects how people frame an interaction, feel, respond and perform, but the effect depends on timing and context, and no single channel carries a fixed meaning.

This context-dependence is most visible when interaction begins. \citet{LavitNicora2024GazeDetection} treat gaze detection as a cue for starting collaboration in an assembly task, where gaze is evidence that a person is ready to coordinate. \citet{AbadMoya2026Initiation} link such detection to legibility. Telling participants even a little about what the robot could perceive made them interpret its engagement responses more consistently, yet 11 of 15 participants still read the robot's orientation plus a short verbal acknowledgement as full engagement, although confirmed engagement came only with a later cue such as a change in LED state. A robot can classify an interaction correctly while still communicating its internal state poorly, so detecting a caregiver and selecting a recovery-oriented action are not sufficient. The robot's visible behaviour must make the energy request, and its urgency, clear on its own.

Time adds a different complication. From the Karotz Home Study, De Graaf, Ben Allouch, and van Dijk \citeyearpar{DeGraaf2016LongTermAcceptance} derive six phases of domestic acceptance, each ending in the user's decision to continue or stop, as the factors driving early adoption give way to others once novelty fades. \citet{Matheus2025LongTerm} review this literature and argue for repeated evaluation in real settings over one-session impressions. An effect measured in a single session may be novelty alone.

Social embeddedness also has to survive absence. \citet{Marquardt2009SituatedMessages} show that situated messages preserve spatial and contextual meaning even when human and robot are not co-present. A recovery architecture needs such a channel when no caregiver is physically available. A notification of battery percentage informs its reader; only a request that ties the robot's state to its consequence and to the help expected tells the reader what to do.

Legibility therefore depends on four things at once: what the robot perceives, when it acts, what its body can show and what people expect of it. The open problem is how to map changes in internal state reliably onto social acts that people can see and that arrive at the right moment.

\section{Cognitive architectures for embodied and socially embedded agents}
\label{sec:soa-architectures}

Modular HRI systems tend to postpone a hard integration problem: how perception, memory, motivation, learning, planning, and motor execution become parts of one continuing agent instead of stages in a pipeline of independent modules. Cognitive architectures address this problem by specifying how those functions constrain one another within an integrated control system \citep{MoulinFrier2018DACH3,GonzalezSantamarta2025Hybrid}. For a socially embedded robot, the question is which architecture stays responsive while still supporting memory, internal valuation, anticipatory action and legible behaviour when interaction is uncertain.

Embodied and enactive approaches reject the idea that cognition can be specified apart from a body and its history. \citet{GeorgeonRitter2012Schema} model emergent cognition through intrinsically motivated schemas whose meaning is constructed through interaction. This position takes on \citeauthor{Harnad1990SymbolGrounding}'s (\citeyear{Harnad1990SymbolGrounding}) symbol-grounding problem: a formal system cannot acquire intrinsic meaning by circulating symbols only among other symbols, so robotic concepts must ultimately be anchored in sensorimotor regularities, outcomes, and the agent's own organisation. \citet{DiPaolo2003OrganismicallyInspired} extends the logic, arguing that adaptive, homeostatic organisation can ground intrinsic teleology, goals that belong to the agent and not to its designer, beyond a closed sensorimotor loop. On this reading, value is what makes a state or action matter to the agent at all, so it cannot be a layer added after perception.

Hybrid robotic architectures put this integration into practice. MERLIN2 \citep{GonzalezSantamarta2025Hybrid} generates, controls, plans and monitors behaviour by coordinating reactive, deliberative and emergent functions in one hybrid architecture, so planning is not the sole source of intelligent conduct. \citet{Fusaro2021HumanAwareBT} take a narrower but complementary route, modifying behaviour trees for human-aware cooperative planning by adding action costs and cost functions for the worker's availability, decisions and ergonomics. Behaviour trees are modular, easy to inspect and predictable in execution. Unless their transitions are linked to a valuation mechanism, their priorities stay where the designer set them.

DAC-h3 is built for proactivity from the outset. \citet{MoulinFrier2018DACH3} integrate reactive control, adaptive learning, contextual knowledge and autobiographical representations in a layered Distributed Adaptive Control architecture that acquires and expresses knowledge about both world and self, so that proactive behaviour can arise from accumulated experience without an immediate trigger. DAC-h3 comes close to the required action-selection loop, but it is domain-general: it leaves open how a single energy-related drive should acquire salience, compete with social opportunities, and produce a socially legible request on iCub specifically.

\citet{Berto2025CuriosityAffect} make motivational structure testable in experiments. Their iCub architecture links sensing and perception to an intrinsic value system built from three drives (survival or recharging, affection or social comfort, and exploration or boredom reduction). By reweighting these drives they produced distinct Playful and Social robot profiles. Participants adjusted their own behaviour and attributed different traits to each profile without being told which was which. The result shows that value-system parameters can reorganise an interaction enough for human partners to infer distinct needs from it. The design also has clear limits: the drive set is small and hand-designed, and scalability, drive conflict, long-term learning, and safety under misclassification remain open.

A different lever comes from \citet{Carminatti2025Stress}, whose stress-inspired mechanisms on iCub alter deliberative functions (attention, memory, decision-making, cognitive flexibility) and leave goal selection itself untouched. For an energy-critical robot such modulation needs limits, since degrading deliberation too aggressively would make the robot least socially appropriate exactly when clear help-seeking matters most.

Emotion architectures joined drives, perception and expression earlier, under a different vocabulary. Kismet \citep{Breazeal2003Emotion} connects affect recognition, drives, emotion dynamics, and expressive behaviour into one loop. \citeauthor{RodriguezRamos2015Emotions}'s (\citeyear{RodriguezRamos2015Emotions}) review of computational emotion models finds the field still short on theoretical grounding, consistent operating cycles, tight coupling with cognition, cross-platform comparison, and validation. \citet{Moerland2018EmotionRL} show that emotion-like variables in reinforcement learning can shape exploration, risk, memory, and action selection. Together, these strands show that affective variables do computational work beyond making a robot look lively to an observer. They also warn against collapsing reward, valence, drive and expression into a single number.

Explainability adds a further architectural requirement. \citet{Stange2022SelfExplaining} propose an explainable behaviour-generation architecture in which a social robot exposes its reasons for acting. A caregiver-mediated recovery system needs two kinds of explanation: behaviour people can read at once through posture, gaze or speech, and, when that reading is uncertain, a plain statement of what the robot lacks and what will happen without help. The value system and action-selection loop must therefore keep their state and priorities traceable. None of these architectures settles how continual context recognition should hand over to value-guided action.

\section{Attention, salience, and engagement selection}
\label{sec:soa-attention}

Attention is how a robot reduces a high-dimensional stream of sensory events to a small set of candidates for action. In social HRI this is unusually hard, because what stands out perceptually is often not what matters for behaviour. A loud sound can be salient without being socially directed; a quiet caregiver entering the room can be the most important event in the scene while motion, speech and task-related stimuli are more salient.

Gaze is the external cue studied most in social HRI, and the work reviewed in Section~\ref{sec:soa-legibility} shows that it cannot be read apart from proxemics and context. As an input to the robot, gaze shows that social attention is more than face detection: the system must estimate whether a person is orienting toward the robot and whether context supports engagement at all. The robot's own kinematics limit how that estimate is formed and shown, which is why \citet{Shaghaghi2024SocialGazeSpace} adapt a Social Gaze Space model specifically to iCub.

More recent work tackles stimulus competition directly. \citet{GarciaMartinez2025BioinspiredStimulus} use reinforcement learning for bio-inspired stimulus selection under multisensory overload in a social companion robot, replacing a fixed priority list with learned selection of whichever input should guide behaviour at that moment. \citet{Ghamati2025LearningToGaze} introduce a habituation-inspired gaze strategy paired with multi-objective Q-learning, so repeated or uninformative stimuli lose priority over time while socially relevant changes regain it. Both gain flexibility at a price: they depend on how the reward is specified and on training conditions that may not match the open-ended statistics of a real environment. A policy trained purely to maximise engagement can also conflict with self-maintenance just when energy is low, a conflict this literature does not address.

The always-on architecture of \citet{Pasquali2025AlwaysOn} takes the opposite approach: its lightweight, non-adaptive front end reduces the scene to five aggregate features that feed a fixed gaze policy (Section~\ref{subsec:intro-modular-organization}). That keeps perception running cheaply, but salience stays almost entirely perceptual. The system can flag that something unusual happened without knowing whether the event is valuable, threatening, socially relevant, or useful for energy recovery. Section~\ref{sec:icub-ecosystem} returns to this gap, which calls for both an action-selection layer and a value system.

None of the systems above predicts whether an approaching person means to interact at all, the task the PAR-D dataset is built for \citep{Thompson2024PARD}. \citet{AbadMoya2026Initiation} show that people draw on several modalities at once to attract a robot's attention, adapting their strategy once the system's detectable cues are disclosed. \citet{Raggioli2025ApproachingBehavior} use deep reinforcement learning to learn approaching behaviour from user activity and disengagement, treating social approach as a policy that must account for its likely effect on the person approached. Salience therefore cannot be determined from the robot's side alone: an interaction is worth starting when a person is available, receptive and well placed. Yet the robot may still need to act when those conditions are only partly met, because its deficit becomes urgent regardless of human readiness.

Salience thus has two sources. External social cues indicate who might be able and willing to help, while internal deviation indicates why that opportunity matters now. The approach and gaze studies reviewed above score candidates from external cues alone, with no input from the robot's own need.

\section{Adaptation, memory, mutual learning, and personalisation}
\label{sec:soa-adaptation}

Adaptation is invoked constantly in HRI and defined precisely far less often. \citet{Ahmad2017Adaptivity} survey systems that adapt to user emotion, personality, behaviour, or interaction history across application domains. They find that most of the reviewed studies rest on short-term interactions, that no common protocol exists for evaluating adaptive robots, and that memory-based adaptation remains under-studied. In this thesis it means something narrower and testable: the robot's stored record of past encounters must reflect participant behaviour and must measurably change a later engagement decision.

Whether that requirement can be met depends on memory architecture. McClelland, McNaughton, and O'Reilly \citeyearpar{McClelland1995ComplementaryLearning} distinguish complementary learning systems, rapid episodic acquisition on one side and slower integration of regularities on the other. The distinction maps directly onto social robots: a single failed recharge interaction should be stored quickly, but a stable belief that one participant is usually responsive should emerge only after repeated evidence. \citeauthor{Pasquali2024ArtificialCognition} (\citeyear{Pasquali2024ArtificialCognition}, \citeyear{Pasquali2025AlwaysOn}) follow the same split for always-on experience: episodes are stored as they close, and clustering consolidates them later.

Mutual-learning studies, all run on a simulated Reachy robot, show how a partner's behaviour enters what the robot learns. \citet{Lemhaouri2022CaregiverResponsiveness} find that a caregiver's contingent responses shape affect-grounded language learning. \citet{Markelius2023MutualLearning} and \citet{Heikkinen2024MutualLearning} combine affect-grounded linguistic interaction with Differential Outcomes Training, in which each need, once met, earns its own distinct feedback. Across these studies the partner's timing, affect and consistency shape what the robot learns, so the partner is more than a source of labels. For caregiver-mediated energy recovery the precedent is direct, since assistance is an outcome of social coordination, and the robot's future behaviour should reflect how particular people have responded before.

Personalisation can target spatial, linguistic, behavioural, or preference parameters, and the literature spreads across all four. \citet{Patompak2020Proxemics} learn personalised proxemics, since acceptable interpersonal distance varies with culture, traits, and familiarity. \citet{PekarekRosin2025DiverseUsers} adapt HRI across diverse user groups by coordinating speech recognition, dialogue, and planning, and \citet{Tanevska2020SociallyAdaptable} build social adaptation into a cognitive framework for iCub (Section~\ref{sec:icub-ecosystem}). \citet{Wang2025PreferenceAction} reduce how much feedback preference-based reinforcement learning needs by learning action representations. These methods differ in machinery but share two conditions: adaptation pays off when the inferred preference is stable enough to justify changing behaviour and the system can fall back to default behaviour if its preference model proves wrong.

Learning from engagement and disengagement bears directly on proactive help-seeking. \citet{Mihhailov2025GuideUnexpected} find that users provide more instruction when a robot's behaviour violates their existing model of it. Read with the approach learning reviewed above \citep{Raggioli2025ApproachingBehavior} and with in situ teaching of social autonomy \citep{Senft2019SocialAutonomy}, the finding suggests that correction, avoidance, delayed assistance, and successful support should enter memory as distinct signals, since a binary success flag would merge them.

Long-term studies explain why memory has to be evaluated across sessions. Beyond the acceptance studies reviewed in Section~\ref{sec:soa-legibility} \citep{DeGraaf2016LongTermAcceptance,Matheus2025LongTerm}, the scoping review of \citet{Newman2025EmpathyLongTerm} shows that empathy-related interaction itself changes across repeated encounters. A policy that looks personalised within one session may only be reacting to cues available in that moment; personalisation is shown only when stored past interaction still changes later behaviour once the original cues are gone.

Adaptive preferences \citep{MarotoGomez2025RobotCompanion} and stable motivational profiles \citep{Berto2025CuriosityAffect} can make a robot's history perceptible, but they describe personality, a property of the robot's parameters. Person-specific memory is a relation between the robot's experience and a particular partner.

Language models can maintain conversational continuity by summarising interaction history, but unconstrained text memory brings risks for privacy and factual accuracy \citep{Wang2026LLMHRISystematic,PintoBernal2025Designing,Ghamati2025PersonalisedLLM}, and in a room with several people it can attach a remark to the wrong person. In a regulatory architecture, generative models can phrase what the robot says but should not decide who helped or whether replenishment succeeded.

Three requirements follow for memory-based adaptation. The stored representation must be compact and outcome-grounded, built from observable response and recovery variables, with no inferred sentiment. Its influence must be testable, so that later engagement differs measurably when the stored history differs. Adaptation must also stay subordinate to safety: the robot may prefer a historically responsive caregiver, yet it must escalate to other people or channels once energy becomes critical.

\section[Multimodal grounding and foundation-model-based HRI]{Multimodal grounding and\\foundation-model-based HRI}
\label{sec:multimodal-grounding}

Language is useful to an embodied robot only once it connects to perception, action and consequence, which is the symbol-grounding problem of Section~\ref{sec:soa-architectures} \citep{Harnad1990SymbolGrounding}. In HRI, grounding means three separate alignments. Words must refer to entities and events actually present in the environment. Plans must be limited to actions the robot can execute. Conversational claims about internal state must correspond to measured variables and controller decisions rather than to whatever a language model finds plausible to say next.

\citet{Allgeuer2024Chatty} tackle these alignments with a modular architecture that grounds an LLM in what the robot perceives and can do. It combines a chat manager with open-vocabulary object detection, pose estimation, gesture detection and other deep models to support open-ended conversation and collaboration. The modularity lets perceptual evidence and action interfaces constrain what the language model may generate. It does not remove the cost of integration: the architecture still coordinates models with different latency, uncertainty, and failure characteristics, so conversational competence is no guarantee of reliable execution under time-critical conditions.

LaMI takes a related route, splitting the problem into scene narration, planning, and expression \citep{Wang2024LaMI}. Separating what the robot perceives from what it intends to do and from how it communicates reduces the burden on any single prompt. \citet{Kim2026Streamlining} integrate LLM-based planning into modular robotic frameworks, and \citet{Kim2024LLMHRI} analyse how LLM-powered HRI changes dialogue, reasoning, task planning, and user experience. This work converges on placing foundation models inside a constrained architecture.

The reviews identify consistent limits. \citet{Wang2026HumanModalities} names multimodal integration, grounding, alignment, and real-time constraints as the central open challenges. \citet{Wang2026LLMHRISystematic} describe a field that is technically broad but uneven in method: user studies differ simultaneously in robot, task, prompting method, model version, and outcome measure, so reported gains in naturalness or task success rarely compare cleanly across papers. \citet{PintoBernal2025Designing} document design challenges around latency, reliability, and embodied context that do not disappear with a better model. \citet{Mauliana2025MultiSession} extend evaluation across sessions, where conversational memory and expectation matter more than in a single encounter. \citet{Ghamati2025PersonalisedLLM} propose a personalised LLM agent for HRI. This work matters for remote communication and for voicing need naturally, but the critical control signal cannot depend on how fluently a model generates text. The architecture must still be able to issue a fixed, verified recovery request when language generation fails.

ROMA streams audio, video and text in real time for interactive understanding and proactive monitoring \citep{Tian2026ROMA}, and VideoMAE learns data-efficient self-supervised video representations for activity recognition \citep{Tong2022VideoMAE}. Such models could sharpen recognition of caregiver availability, interaction onset, or response behaviour, but their value depends on computational cost, inference latency, hardware requirements, and robustness to distribution shift. \citet{Pasquali2025AlwaysOn} choose lightweight perceptual features for these practical reasons: an always-on robot cannot assume unlimited compute.

Multimodal communication covers affect and non-verbal expression as well as vision and language. In the mutual-learning studies of Section~\ref{sec:soa-adaptation}, learning on both sides improved when each satisfied need earned its own sound and movement instead of one shared positive response \citep{Markelius2023MutualLearning,Heikkinen2024MutualLearning}, because the extra channel carried need-specific information. For expressing internal state, the matching requirement is that the channels agree: verbal urgency, gaze direction, posture, facial expression, and remote notification must signal the same state, or the redundancy turns into noise.

Remote communication faces the grounding problem of the situated messages discussed in Section~\ref{sec:soa-legibility} \citep{Marquardt2009SituatedMessages}. Foundation models can improve a remote channel by tailoring wording, summarising context, or choosing a suitable modality. They should not be allowed to invent the internal condition, suppress a critical alert, or select an unverified physical action. The division of labour is deliberately asymmetric: deterministic modules own measurement, thresholds, safety and action permissions, learned multimodal modules estimate social context, and generative models word the robot's messages within bounds the other two have set. Still unresolved is how these three kinds of module should exchange state with memory, the body and the running control loop.

\section{Homeostasis, allostasis, and drive-based artificial motivation}
\label{sec:soa-homeostasis}

Attention, memory and language give a robot the machinery of social action without giving it a reason to act. Physiology confronted the question of where such reasons come from long before robotics did. \citet{Cannon1929Homeostasis} described the coordinated organisation through which open biological systems keep internal variables within viable ranges. The formulation is often reduced to fixed-point error correction, though even classical homeostasis concerns dynamic stability and never implied that internal variables stay fixed. \citeauthor{Tinbergen1951Instinct}'s (\citeyear{Tinbergen1951Instinct}) ethological account adds the missing behavioural mechanism: internal motivational states alter the readiness of action systems, while external releasing stimuli determine when and how behaviour is expressed. This pairing of internal deficit with outside opportunity carried over into social robotics: Kismet's drives and triggering events draw on the same ethological models \citep{Breazeal2003Emotion}.

Allostasis keeps the vocabulary but changes the timing. \citet{Sterling2012Allostasis} argues that efficient regulation is predictive: instead of waiting for an error and then restoring a constant set point, the organism anticipates demand, reallocates resources in advance, and adjusts its own operating ranges. \citet{SchulkinSterling2019Allostasis} develop the brain-centred version, linking prediction, prioritisation, learning, and dopaminergic reinforcement. \citet{Hulme2019HomeostaticControl} survey the wider neurocomputational literature, in which internal-state control ranges from reactive feedback to model-based and predictive policies. \citet{Petzschner2021Interoception} draw a parallel distinction between reflexive control and flexible regulation informed by interoceptive context, experience, and expectation. For robotics, a homeostatic controller acts because a measured variable has already departed from a preferred range, while an allostatic controller acts because it predicts that inaction will produce a costly deviation later.

The free-energy and active-inference literature generalises the same idea with generative models of internal and external causes. It relates physiological regulation to variational dynamics \citep{BettingerFriston2023PhysiologicalRegulation}, formalises how a system maintains its characteristic states \citep{Friston2023FreeEnergy}, and presents active inference as a single account of perception and action \citep{Pezzulo2024ActiveInference}. These accounts give a principled way to combine uncertainty, prediction and preference, though empirical support varies sharply across specific predictive-processing claims \citep{Hodson2024PredictiveCoding}. An architecture built around a single energy set point gains nothing from invoking active inference as a label: it would have to implement the generative model and test the predictions that follow, or the vocabulary would claim more than the architecture does.

Robotic homeostatic controllers offer more direct precedent. \citet{Cao2017CollaborativeHomeostatic} combine homeostatic variables with reactive and deliberative layers in a collaborative social-robot controller, and a related end-user interface lets non-programmers configure homeostatic behaviour for NAO in robot-assisted social therapy \citep{Cao2017EndUserHomeostatic}. Both show that internal variables can drive socially observable action, and that authoring such behaviour need not require touching low-level code. Their shared limitation is that structure, thresholds, and behavioural repertoire are largely preconfigured: neither offers a general method for learning which social action will restore a given internal resource.

\citet{GuerreroRosado2021RegulatoryBehaviour} propose a biologically inspired regulatory architecture grounded in the medial reticular formation and implement an initial robot model with simulated hunger and thirst. \citet{Ngo2021DesertReptiles} tests homeostatic and allostatic principles in a robotic model of desert-reptile regulation, asking whether the controller reproduces the species' typical behaviour under environmental constraints. \citet{Rosado2022DriveCompetition} focus on competition among drives, using mutually inhibitory neural populations to arbitrate efficiently among needs; their analysis of efficiency, fairness, and stability matters once more than one need can dominate action. A single orexigenic drive avoids instability between competing needs, but it still competes with ongoing social interaction, task commitments, and safety constraints, so its priority has to be graded.

For \citet{DiPaolo2003OrganismicallyInspired}, homeostasis is definitional: an agent is meaningfully adaptive only when it judges events by whether they preserve its own organisation. \citet{Pasquali2026ISEED} arrive at the same requirement from the direction of cognitive architecture: without an evaluative layer, a system can predict what may happen but cannot decide what should be pursued. Regulatory variables are attractive here because their state belongs to the agent, even when a designer has fixed the set point.

HRI research has also built artificial drives and related affective variables into social robots. \citet{LewisCanamero2016Hedonic} separate hedonic quality from externally specified reward in a motivated robot (Section~\ref{sec:soa-expression}). \citet{MarotoGomez2018MotivationalDecision} propose a bio-inspired motivational decision system in which user perception feeds back into the robot's own motivational state. Later work from the same group adds biologically inspired decision-making, artificial biological functions, adaptive circadian rhythms, dopamine-like modulation, and adaptive object preferences with emotional responses \citep{MarotoGomez2023BiologicalDecision,MarotoGomez2023ArtificialBiologicalFunctions,MarotoGomez2023Circadian,MarotoGomez2024Dopamine,MarotoGomez2025RobotCompanion}.

\citet{Yoshida2025HomeostasisRL} link need to learned behaviour formally through homeostatically regulated reinforcement learning (HRRL), where reward is defined by drive reduction: an outcome is valuable when it moves the internal state toward its preferred region, aversive when it increases deviation. Under the paper's stated assumptions, learned policies minimise discounted future deviation. This ties drive reduction to optimal feedback control and predicts state-dependent value, risk aversion, competition between needs, and anticipatory regulation wherever predictive cues are available.

HRRL is a strong normative reference, but it settles fewer design choices than it appears to. Drive function, set point, discounting, state representation, action space, and available resources must still be specified by hand, and \citet{Yoshida2025HomeostasisRL} flag open questions of their own: storable resources, event-driven consumption, multiple timescales, biological realism, integration with active inference. An energy-centred approach has an advantage here: the energy variable is directly measurable and the replenishment outcome concrete, so drive reduction can be read directly from the log. The remaining difficulty is coordinating social recovery behaviour under a caregiver's uncertain availability.

\citet{Scheunemann2022IntrinsicallyMotivated} add that participants rated intrinsically motivated autonomous behaviour significantly warmer than a reactive baseline, suggesting that behaviour driven from within can make a robot seem more like a social other. But intrinsic motivation is not automatically a benefit. Behaviour that originates inside the robot is also harder for a person to predict, whereas scripted behaviour can be made fully legible. The design target is therefore constrained autonomy: behaviour that originates in the robot's internal regulatory state while staying interpretable, safe, and responsive to the person in front of it.

The social environment adds norms that most resource-foraging models never face. A robot can have a strong energy deficit while the person nearby is engaged in a sensitive conversation, physically distant, or otherwise unavailable. \citet{Senft2019SocialAutonomy} address a version of this tension by letting people teach social autonomy in situ. This suggests a practical route: caregiver responses can shape the timing or choice of recipient without touching the hard safety requirement that a critical deficit must eventually be communicated to someone.

Together this work supplies the basis for an orexigenic drive. Classical homeostasis contributes monitoring, deficit detection and corrective action, allostasis adds anticipation, drive competition explains selection under conflict, and HRRL shows how internal improvement can define reward. None of it delivers the social recovery loop itself: an embodied robot that converts a measurable energy deficit into a proportionate request for human help, learns from how people respond to that request, and keeps a transparent link between internal state and outward conduct throughout.

\section{Social expression and explanation of internal robot states}
\label{sec:soa-expression}

An internal variable matters in interaction only when its changes show up as behaviour a person can read and act on. That mapping decides whether a human can distinguish ordinary variation from a condition needing assistance, whether the robot's request is proportionate to its urgency, and whether the resulting response closes the regulatory loop. Kismet set the lasting pattern \citep{Breazeal2003Emotion}: facial and vocal behaviour there regulated the interaction by shaping caregiver response, which made expression part of control. Pre-authored displays \citep{Cao2017CollaborativeHomeostatic,Cao2017EndUserHomeostatic} show the cost of the opposite choice: a display tied to a fixed threshold ignores interaction history and whether a caregiver is available.

Recharging research offers the most direct precedent for energy-centred recovery. After measuring how recharging degraded a robot's social service over three weeks in an office, \citet{Deshmukh2018Recharging} tested verbal strategies (transparency, apology, politeness) for managing user expectations during recharge, and these made the robot more acceptable while it charged. The work also reframes recharging as a visible social event that affects expectations, relationships and perceived competence.

What should be expressed follows from the distinction between reward and hedonic quality. \citet{LewisCanamero2016Hedonic} model basic pleasure as the direction and magnitude of homeostatic change, independent of any fixed outcome label, so the same event can be beneficial or harmful depending on current state. What is worth communicating is how far the robot is from its set point, where it is heading, and which action would reduce the deficit. A flat report that the battery is low conveys none of these. Because attention and confirmed engagement are distinct states, explanation has to back up expression wherever cues remain ambiguous \citep{AbadMoya2026Initiation,Stange2022SelfExplaining}.

Empathy research draws an ethical boundary here as much as an interactional one \citep{Park2022Empathy,Newman2025EmpathyLongTerm}. A recovery-oriented robot can behave in a socially sensitive way without simulating suffering or manipulating guilt. The target is a request for care proportioned to the need, where the user understands the need and the consequence of inaction while keeping the freedom to decline.

Two outcomes must then be kept apart. Successful expression should help people notice the need, respond appropriately and bring about verified recovery, without excessive interruption or coercive urgency. Whether it also makes the robot seem natural is a separate question. Perception ratings such as the warmth reported by \citet{Scheunemann2022IntrinsicallyMotivated} and the naturalness reported by \citet{MarotoGomez2025RobotCompanion} can reflect coherence, variability, anthropomorphic attribution or novelty. They do not reveal which internal mechanism participants understood, nor whether the behaviour obtained anything the robot needed.

\section{The iCub ecosystem for cognitive and interaction architectures}
\label{sec:icub-ecosystem}

iCub was designed around a specific research programme: cognition should be studied through development, sensorimotor experience and interaction on an open humanoid platform. Sandini, Metta, and Vernon \citeyearpar{Sandini2007ICubCognitive} present iCub as an open system for enactive cognition, and \citet{Metta2010ICub} describe the mature platform as an open-systems testbed for cognitive development. Humanoid morphology, an articulated head and eyes, upper-body expressivity, tactile sensing, and the YARP middleware together make iCub suited to work in which perception, action, and social interpretation must be studied jointly.

The same features also limit the present results. Gaze is legible only through coordinated eye and head movement, so the direction of attention is harder to read than on a human face. Onboard computation is limited relative to contemporary foundation models, so processing is spread across external machines, at the cost of added latency and dependence on infrastructure. The humanoid appearance raises user expectations that can far exceed actual autonomy. Rich embodiment makes iCub results valuable, yet they cannot be assumed to generalise to mobile service robots with different sensors, mobility, or expressive channels.

\citet{Fischer2018ICubHRI} address integration through iCub-HRI, a software framework for complex HRI scenarios that reduces the cost of combining perception, dialogue, and behaviour through reusable modules and interfaces. The contribution is infrastructural: complex scenarios become implementable, while what the robot should value during them is left open.

In the socially adaptable framework of \citet{Tanevska2020SociallyAdaptable}, an internal comfort level driven by the caretaker's stimuli guides iCub's social adaptation, and user studies compare adaptive with non-adaptive behaviour. Closest in spirit to a regulatory variable is the comfort drive of \citet{Mongile2025ComfortDriven}, derived from attachment theory, which registers how pleasant iCub finds an interaction and steers its adaptation to the partner's interaction style. Users read the resulting behaviour as more goal-oriented and affective. With the drive profiles of \citet{Berto2025CuriosityAffect} and DAC-h3 \citep{MoulinFrier2018DACH3}, these frameworks show that iCub can carry the loop from internal variables to socially perceived action, though each implements a different part of it and evaluates it differently.

The Artificial Cognition programme of \citeauthor{Pasquali2024ArtificialCognition} (\citeyear{Pasquali2024ArtificialCognition}, \citeyear{Pasquali2025AlwaysOn}) is the architecture this work builds on directly; Section~\ref{sec:intro-developmental-architecture} describes its implementation. Its strengths are sustained operation, situated first-person data, and self-supervised organisation of experience grounded in the robot's personal history. Its limitations define the present contribution. The system has no explicit goal beyond detecting novelty and regularity. Segmentation fragments high-activity periods. Memory has no separate short- and long-term dynamics and grows with no forgetting, reinforcement or value to decide which episodes to keep. Semantic information stays limited, and recognition never connects to action.

\citet{Pasquali2026ISEED} consolidate the programme as iSEED, promoting values to a pillar alongside embodiment and memory on the grounds that time, body and memory organise experience but do not decide how the agent behaves. Measured against their criterion of causal grounding, the face-weighted segmentation is only a proto-value system, sufficient to bias what gets remembered and no more. Section~\ref{sec:intro-homeostasis-drive} sets out how the orexigenic drive answers that criterion.

\citet{Sandini2025MutualUnderstanding} situate these developments within a broader programme of shared embodied cognition, arguing for architectures combining first-person experience, prospection, episodic memory, and social interaction. \citet{Morasso2026BioInspiredCognitive} contrasts bio-inspired cognitive robotics with embodied AI, favouring socially acceptable, resource-aware intelligence over disembodied scale. Both support extending an always-on embodied architecture rather than moving all control inside a foundation model.
